\documentclass[12pt]{article}
\usepackage{amsmath}

\usepackage{xcolor}

\usepackage{listings}
\usepackage[most]{tcolorbox}

% style stolen from haskell notes (non avrei ricreato appositamente altrimenti...)
\definecolor{mzbg}{HTML}{F8F9FA}
\definecolor{mzpurple}{HTML}{5E5086}
\definecolor{mzgreen}{HTML}{118C4F}
\definecolor{mzred}{HTML}{C0392B}
\definecolor{mzblue}{HTML}{2E86C1} 
\definecolor{mznumber}{HTML}{ADB5BD}
\definecolor{mztitlebg}{HTML}{2C3E50}

\lstdefinelanguage{MiniZinc}{
    morekeywords={include, var, par, array, of, int, float, set, bool, constraint, solve, satisfy, minimize, maximize, forall, exists, where, if, then, else, endif},
    morekeywords=[2]{alldifferent, abs, sum, count, min, max, card}, 
    sensitive=true,
    morecomment=[l]{\%},
    morecomment=[s]{/*}{*/},
    morestring=[b]",
}

\lstdefinestyle{mznstyle}{
    language=MiniZinc,
    commentstyle=\color{mzgreen}\itshape,
    keywordstyle=\color{mzpurple}\bfseries,
    keywordstyle=[2]\color{mzblue},
    numberstyle=\tiny\color{mznumber},
    stringstyle=\color{mzred},
    basicstyle=\ttfamily\footnotesize,
    breaklines=true,
    showstringspaces=false,
    tabsize=2,
    numbers=none
}

\newtcblisting{mznbox}[1]{
    enhanced,
    skin=bicolor,
    colback=mzbg,
    colbacktitle=mztitlebg,
    coltitle=white,
    colframe=mztitlebg,
    boxrule=0.5pt,
    arc=4pt,
    drop shadow=black!10!white,
    fonttitle=\bfseries\sffamily,
    title=#1,
    listing only,
    listing options={style=mznstyle, numbers=left, numbersep=10pt},
    left=15pt,
    top=4pt,
    bottom=4pt,
    breakable
}

\title{E.A.5.6 - CSP: Cards}
\author{}
\date{}

\begin{document}
\maketitle
\vspace{-5em}
\section{Modellazione}

\begin{itemize}
    \item \textbf{variabili}: $X = \{X_{i,j} \mid 1 \leq i \leq n, 1 \leq j \leq k\}$ - ogni variabile $X_{i, j}$ rappresenta la $j$-esima copia della carta con il numero $i$

    \item \textbf{domini}: $D_{X_{i,j}} = \{1 \dots k \times n\}$ per tutte le variabili $X_{i, j}$ - rappresentano le posizioni dove si possono collocare le carte

    \item \textbf{vincoli}: 
        \begin{itemize}
            \item alldiff($X_{1,1}, \dots X_{n, k}$) (in ogni posizione ci può essere al massimo una carta)
            \item $X_{i, j+1} - X_{i,j} = i+1$ \ per $1 \leq i \leq n, 1 \leq j \leq k-1$ 
                \subitem vincolo sulla distanza, in cui introduciamo anche una concezione di ordine che ci fa da ``symmetry-breaker'': le carte sono disposte in ordine di copie
        \end{itemize}
\end{itemize}


\section{Istanziazione}

\begin{itemize}
    \item \textbf{variabili}: \\
     $n=4$ carte diverse e $k=2$ copie per ciascuna, $2 \times 4 = 8$ variabili:
    $$X = \{X_{1,1}, X_{1,2}, X_{2,1}, X_{2,2}, X_{3,1}, X_{3,2}, X_{4,1}, X_{4,2}\}$$
    
    \item \textbf{domini}: \\
        $$D_{X_{i,j}} = \{1, 2, 3, 4, 5, 6, 7, 8\} \quad \text{per ogni } X_{i,j} \in X$$
    
    \item \textbf{vincoli}:
    \begin{itemize}
        \item posizioni
            $$\text{alldiff}(X_{1,1}, X_{1,2}, X_{2,1}, X_{2,2}, X_{3,1}, X_{3,2}, X_{4,1}, X_{4,2})$$
        
        \item distanza e simmetrie:
        \begin{align*}
            X_{1,2} - X_{1,1} &= 2 \quad \text{(1 carta a dividere le due copie di 1)} \\
            X_{2,2} - X_{2,1} &= 3 \quad \text{(2 carte a dividere le due copie di 2)} \\
            X_{3,2} - X_{3,1} &= 4 \quad \text{(3 carte a dividere le due copie di 3)} \\
            X_{4,2} - X_{4,1} &= 5 \quad \text{(4 carte a dividere le due copie di 4)}
        \end{align*}
    \end{itemize}
\end{itemize}

\section{Codifica in MiniZinc}

\begin{mznbox}{}
include "alldifferent.mzn";

int: n;
int: k;

array[1..n, 1..k] of var 1..(n*k): X;

% ho trovato array1d nei built-in (rende una matrice un array monodimensionale)
constraint alldifferent(array1d(X));

constraint forall(i in 1..n, j in 1..k-1)(
    X[i, j+1] - X[i, j] = i+1
);

solve satisfy;
\end{mznbox}

\end{document}
